\scanpage{156}
\uline{\unsure{Duals|Duels} et Homs de motifs}

On suppose que\inserted{ $\rho=2$ et que} pour tout motif effectif $M$, de degré $i\geq0$, $\mathcal Hom(M,\Lambda(-i))$ existe\inserted{ et est un motif eff.} (où $\mathcal Hom(M,N)$ est défini axiomatiquement par
\[
\operatorname{Hom}(P,\mathcal Hom(M,N))\simeq\operatorname{Hom}(P\otimes M,N)
\]
On constate que si $M,N$ sont des motifs effectifs de degrés $i,j$, tels que $\mathcal Hom(M,N)$ existe, alors $\mathcal Hom(M,N)$ est de degré $j-i$. Ici, posant
\[
M^{*}=\mathcal Hom(M,\Lambda(-i))\qquad\deg M=i
\]
on trouve que $\deg M^{*}=\deg M$. On trouve d'ailleurs que si $\mathcal Hom(M,N)$ existe dans $\mathcal M$, il existe aussi dans $\widehat{\mathcal M}$ et est le même.

On voit que pour \unsure{tt|$\forall$} $M$ de degré $i$, et \unsure{tt|$\forall$} $N$, $\mathcal Hom(M,N(-i))$ existe, car on aura
\[
\mathcal Hom(M,N(-i))\simeq M^{*}\otimes N
\]
(vérification facile). \deleted{D'où} Cela prouve que pour tout $N$ (pas néc. \unsure{hom.})
\mspage{155}
